\documentclass[11pt,a4paper]{article}

\usepackage{fontspec}
\setmainfont{Times New Roman}

\usepackage{amsmath,amssymb,amsfonts}
\usepackage{geometry}
\geometry{a4paper, margin=0.8in}
\usepackage{booktabs}
\usepackage{tabularx}
\usepackage{array}
\usepackage{listings}
\usepackage{xcolor}
\usepackage{enumitem}
\usepackage{fancyhdr}
\usepackage[colorlinks=true, linkcolor=blue, urlcolor=blue]{hyperref}

\pagestyle{fancy}
\fancyhf{}
\rhead{\small\textit{Báo Cáo Thực Nghiệm Khoa Học}}
\lhead{\small Antigravity Assignment Solver}
\cfoot{\thepage}

\lstset{
    basicstyle=\ttfamily\footnotesize,
    breaklines=true,
    frame=single,
    backgroundcolor=\color{black!3},
    columns=fullflexible
}

\begin{document}

\begin{center}
    {\LARGE\textbf{BÁO CÁO THỰC NGHIỆM VÀ PHÂN TÍCH}} \\[0.5em]
    {\normalsize Môn học: Xử Lý Ngôn Ngữ Tự Nhiên / Machine Learning} \\[0.3em]
    {\small\textit{\today}}
\end{center}

\vspace{0.8em}
\hrule
\vspace{1.2em}

\section*{1. Mục Tiêu và Cơ Sở Lý Thuyết}
Nội dung phân tích lý thuyết, biểu diễn toán học và các công thức liên quan.

\section*{2. Tính Toán Chi Tiết (Hand Calculations)}
Trình bày đầy đủ từng bước thay số và biến đổi đại số.

\section*{3. Kết Quả Thực Nghiệm}
Trình bày bảng số liệu và phân tích metrics.

\section*{4. Phân Tích Lỗi và Bài Học (Error Analysis \& Reflection)}
Đánh giá định tính và đối chiếu lý thuyết.

\end{document}
